\subsection{Summary and Research Gaps}

Table~\ref{tab_sbl_literature_summary} highlights the distinct priority-restoration mechanisms of IBL, BLIP, and DBL. Despite these differences, the reviewed studies leave three related gaps between the operational representation of shared bus lanes and the analytical conditions required for their implementation.

First, the effects of traffic signals, bus-stop dwell time, and the finite duration of lane-changing maneuvers have not been integrated within a unified analytical feasibility framework. Existing studies have incorporated some of these factors separately. For example, signal timing and bus-stop operations have been considered in analytical or optimization models of IBL \citep{lin2022capability,xie2023modelling}, while connected-vehicle studies have developed state-responsive clearance rules \citep{xie2021design,jiang2025analysis}. Their interactions, however, remain insufficiently characterized. Signal timing determines the available travel and clearance windows, bus dwell time changes the effective spacing between consecutive buses, and lane-changing duration determines whether vehicles can leave the priority space before interfering with bus operation. Treating these factors separately may therefore misrepresent the conditions under which lane sharing is operationally feasible.

Second, the delay externality of mandatory lane changes has not been explicitly incorporated into analytical evaluations of shared bus lane operation. Microscopic simulations can reproduce gap acceptance and lane-changing behavior, but they generally report aggregate traffic outcomes rather than express the resulting disturbance as an analytical cost. In particular, a vehicle leaving the shared lane may force vehicles in the receiving general-purpose lane to decelerate, create gaps, or propagate a delay wave. Empirical and trajectory-based studies show that the magnitude and persistence of this disturbance increase with traffic density \citep{laval2006lane,tian2016empirical,yang2025impact}. Without explicitly accounting for this external cost, existing feasibility conditions capture the benefit obtained by lane-sharing vehicles more readily than the delay transferred to general traffic.

Third, existing studies do not provide a closed-form operational threshold based on net system benefit. Analytical work has derived capacity bounds for the number of private vehicles that can use an intermittent bus lane without increasing bus delay \citep{lin2022capability}, while simulation and optimization studies have identified feasible operating ranges under specified traffic conditions \citep{xie2023modelling,wu2017evaluating,othman2023dynamic}. These results do not, however, yield a transferable critical lane-sharing density that simultaneously accounts for the travel-time savings of lane-sharing vehicles and the lane-changing delay imposed on the receiving lane. Such a threshold is needed to determine when additional lane sharing ceases to produce a non-negative system benefit.

The present study addresses these gaps by developing a unified analytical framework for Shared Bus-Only Lanes (SBL). The framework jointly represents signal timing, bus-stop dwell time, and finite lane-changing duration, and explicitly accounts for the delay imposed on the receiving general-purpose lane. Lane sharing is considered beneficial only when the travel-time savings obtained by lane-sharing vehicles are no smaller than the associated delay to general traffic. Setting these benefits and costs equal yields closed-form expressions for the critical lane-sharing density as a function of general-lane traffic density. The model is developed in three stages of increasing operational detail: a basic corridor, an extension with signalized intersections, and a full formulation incorporating bus-stop dwell time.
